\pdfbookmark{Общая характеристика работы}{characteristic}
\section*{ОБЩАЯ ХАРАКТЕРИСТИКА РАБОТЫ}

\newcommand{\actuality}{\pdfbookmark[1]{Актуальность}{actuality}\textbf{{\actualityTXT}}}
\newcommand{\progress}{\pdfbookmark[1]{Разработанность темы}{progress}\textbf{{\progressTXT}}}
\newcommand{\aim}{\pdfbookmark[1]{Цели}{aim}{{\textbf\aimTXT}}}
\newcommand{\tasks}{\pdfbookmark[1]{Задачи}{tasks}\textbf{{\tasksTXT}}}
\newcommand{\novelty}{\pdfbookmark[1]{Научная новизна}{novelty}\textbf{{\noveltyTXT}}}
\newcommand{\influence}{\pdfbookmark[1]{Практическая значимость}{influence}\textbf{{\influenceTXT}}}
\newcommand{\methods}{\pdfbookmark[1]{Методология и методы исследования}{methods}\textbf{{\methodsTXT}}}
\newcommand{\defpositions}{\pdfbookmark[1]{Положения, выносимые на защиту}{defpositions}\textbf{{\defpositionsTXT}}}
\newcommand{\reliability}{\pdfbookmark[1]{Достоверность}{reliability}\textbf{{\reliabilityTXT}}}
\newcommand{\probation}{\pdfbookmark[1]{Апробация}{probation}\textbf{{\probationTXT}}}
\newcommand{\contribution}{\pdfbookmark[1]{Личный вклад}{contribution}\textbf{{\contributionTXT}}}
\newcommand{\publications}{\pdfbookmark[1]{Публикации}{publications}\textbf{{\publicationsTXT}}}

\input{common/characteristic} % общий с введением диссертации файл

\pdfbookmark{Содержание работы}{description}
\section*{СОДЕРЖАНИЕ РАБОТЫ}

\noindent\textbf{Основные результаты работы:}
%% Список основных результатов работы.
%% Подключается в двух местах: Dissertation/conclusion.tex (заключение
%% диссертации) и Synopsis/content.tex (заключение автореферата).
%% Согласно ГОСТ Р 7.0.11-2011, п. 5.3.3 и 9.2.3.
\begin{enumerate}
    \item Установлено, что стандартная специальная обертывающая алгебры Ли
          $N\Phi(\mathbb{K})$ существует для всех типов $\Phi$, кроме типов $D_n$ $(n\geq
          4)$ и $E_n$ $(n=6,7,8)$; для типа $F_4$ построены как стандартные, так и
          нестандартные специальные обертывающие алгебры.

    \item Для алгебр Ли $N\Phi(\mathbb{K})$ классических типов найдены условия
          однозначности специальных обертывающих алгебр $R_{\Phi}$. Специальная
          обертывающая алгебра типа $A_{n-1}$ $(n>3)$ стандартна тогда и только тогда,
          когда она или ее противоположная алгебра по модулю аннулятора изоморфна
          $NT(n,\mathbb{K})/\mathbb{K}e_{n1}$; ассоциативная среди стандартных
          единственна и изоморфна $NT(n,\mathbb{K})$. Для типа $B_n$ среди
          стандартных специальных обертывающих алгебр, у которых
          $R_{\Phi}/T_{2,-1}\simeq NT(n+1,\mathbb{K})$, условие монотонности
          "--- вложение каждого члена последовательности подалгеброй в
          следующий "--- однозначно выделяет алгебру $RB_n(\mathbb{K})$.

    \item Доказано, что алгебра $NT(n+1,\mathbb{K})$ "--- единственная с точностью до
          изоморфизма ассоциативная точная обертывающая алгебры Ли $N\Phi(\mathbb{K})$
          типа $A_n$ $(n\geq 3)$, допускающая градуировку, согласованную с градуировкой
          Картана. Условие стандартности такой алгебры слабее и однозначности уже не
          дает.

    \item Для алгебр Ли $N\Phi(\mathbb{K})$ типов $B_n$, $C_n$ и $D_n$ над произвольным
          полем построены левосимметрические точные обертывающие алгебры в виде
          полупрямых сумм алгебры $NT(n,\mathbb{K})$ и подходящего идеала алгебры Ли
          $N\Phi(\mathbb{K})$.

    \item Решена проблема~1 из~\cite{sigsam2001}: найдено число стандартных идеалов
          алгебры Ли $N\Phi(q)$ над конечным полем для всех классических и исключительных
          типов $\Phi$. Перечислены нестандартные идеалы обертывающей алгебры $RD_n(q)$,
          чем завершено перечисление идеалов специальных обертывающих алгебр $R_\Phi$
          классических типов; при этом впервые вычислена трехкратная комбинаторная сумма
          с $q$-биномиальными коэффициентами методом интегрального представления
          комбинаторных сумм.
\end{enumerate}


Во \textbf{введении} обосновывается актуальность темы, приводится обзор
результатов об идеалах алгебр $N\Phi(\mathbb{K})$ и о нормальных подгруппах
унипотентных подгрупп групп Шевалле, формулируется цель работы, излагаются
научная новизна и практическая значимость, а также вводятся идеалы $T(r)$,
$Q(r)$, $T(\mathcal{L})$, $Q(\mathcal{L})$, понятие множества углов
$\mathcal{L}(H)$ и понятие стандартного идеала.

\medskip

\textbf{Первая глава} посвящена специальным точным обертывающим алгебрам
$R_{\Phi}$ алгебры Ли $N\Phi(\mathbb{K})$ и вопросам их существования и
однозначности.

В \S~1.1 при фиксированном выборе знаков структурных констант $N_{rs}$ вводится
$\mathbb{K}$-алгебра $R_{\Phi}$ с базой $\{e_r \mid r \in \Phi^+\}$ и
умножением $e_re_s=0$ при $r+s\notin\Phi$ и
\[
    e_r e_s = e_{r+s}, \qquad e_s e_r = (1 - N_{rs}) e_{r+s}
    \qquad (r, s, r+s \in \Phi^+,\ N_{rs} \geq 1);
\]
такие алгебры называем \emph{специальными обертывающими}.

\dissth{p:mult-enveloping}

Из следующей леммы видно, что для типов с корнями разной длины среди
специальных обертывающих ассоциативных нет.

\dissth{l:Const-RPhi-Non-assoc}

В \S~1.2 разобран исключительный тип $F_4$: знаки констант $N_{rs}$ на
экстраспециальных парах корней подбираются так, чтобы получить обертывающую с
нестандартными идеалами и, наоборот, без них.

\dissth{th:stand-f4}

В \S~1.3 получены теоремы единственности для классических типов. Различные
корни $r,s\in\Phi^+$ называем \emph{$p$-связанными} для простого корня $p$,
если $r+p,s+p\in\Phi^+$. Основное средство "--- перенос стандартности на
подалгебры $R(\Psi)=\sum_{a\in\Psi^+}\mathbb{K}e_a$, отвечающие подсистемам
$\Psi$ типа $A_3$ с базой из $p$-связанных корней $r,s$ и простого корня $p$:
стандартность такой подалгебры равносильна равенству $N_{s,p}=N_{p,r}$, то есть
тому, что $e_r$ и $e_s$ лежат в разных односторонних аннуляторах элемента
$e_p$.

\dissth{th:Stand-An}

\dissth{p:Exc-Dn-En}

Последовательность обертывающих алгебр $R_n$ одного классического типа
возрастающих рангов $n$ называем \emph{монотонной}, если каждый ее член
изоморфно представляется подалгеброй следующего члена. Для типа $B_n$
рассмотрим условие
\begin{equation*}
    R_{\Phi}/T_{2,-1}\simeq NT(n+1,\mathbb{K}),
    \tag{\ref*{IsomR-Bn}}
\end{equation*}
где $T_{2,-1}\coloneq T(2p+q)$, $p$ "--- короткий простой корень, а $q$ "---
соседний с ним простой корень. Условие монотонности позволяет получить
следующую теорему единственности.

\dissth{th:Stand-Bn}

\medskip

\textbf{Вторая глава} посвящена точным обертывающим алгебрам, согласованным с
градуировкой Картана.

В \S~2.1 изложен аппарат градуированных алгебр. Для абелевой группы $G$
\emph{$G$-градуировкой} алгебры $A$ называем разложение $A=\bigoplus_{g\in
    G}A_g$ с условием $A_gA_h\subset A_{g+h}$. Градуировка $\Gamma_2$
\emph{уточняет} $\Gamma_1$, если есть гомоморфизм огрубления, вкладывающий
компоненты $\Gamma_2$ в компоненты $\Gamma_1$; градуировка \emph{тонкая}, если
собственных уточнений у нее нет, и имеет \emph{ширину один}, если все ее
ненулевые компоненты одномерны. Для интересующего нас класса алгебр эти два
свойства совпадают.

\dissth{p:fine-graded-nilpotent}

Базис Шевалле задает на $N\Phi(\mathbb{K})$ тонкую \emph{градуировку Картана}
$\Gamma_C\colon N\Phi(\mathbb{K})=\bigoplus_{r\in\Phi^+}\mathbb{K}e_r$. Тонкую
градуировку Ли-допустимой алгебры $R$, для которой перенесенная на $R^{(-)}$
градуировка эквивалентна $\Gamma_C$, называем \emph{согласованной с
    градуировкой Картана}.

\dissth{l:enveloping-roots-grading}

В \S~2.2 отсюда выводится основной результат главы; предварительно
устанавливается, что всякая ассоциативная точная обертывающая для типа $A_n$
$(n\ge2)$ нильпотентна.

\dissth{th:Stand-Graded-An}

Условие градуировки здесь существенно: без него ассоциативная точная
обертывающая типа $A_3$, не изоморфная $NT(4,\mathbb{K})$, существует, а при
$n=2$ и $\Char\mathbb{K}\neq2$ имеются две неизоморфные ассоциативные точные
обертывающие с такой градуировкой.

\medskip

\textbf{Третья глава} посвящена левосимметрическим (прелиевым) структурам.

Через $L_x\colon y\mapsto xy$ $(x,y\in R)$ обозначим оператор левого умножения
в алгебре $R$.

\dissdef{def:pre-lie-algebra}

Всякая левосимметрическая алгебра Ли-допустима, а класс левосимметрических
алгебр шире класса ассоциативных. В \S~3.1 вводится полупрямая сумма
левосимметрических алгебр, с помощью которой в \S~3.2 строятся точные
обертывающие алгебры для типов $B_n$, $C_n$ и $D_n$.

\dissdef{def:semidirect-pre-lie}

\dissth{th:semidirect-assoc-cond}

\dissth{th:semidirect-pre-lie}

В \S~3.2 конструкция применяется к классическим типам. Базу системы корней
$\Phi$ типа $D_n$, $C_n$ или $B_n$ записываем в виде
$\Pi(\Phi)=\Pi(A_{n-1})\cup\{r_\Phi\}$, где
\[
    r_\Phi\coloneq
    \begin{cases}
        \epsilon_2 + \epsilon_1, & \Phi = D_n, \\
        2\epsilon_1,             & \Phi = C_n, \\
        \epsilon_1,              & \Phi = B_n,
    \end{cases}
\]
и получаем разложение $N\Phi(\mathbb{K})\simeq NT(n,\mathbb{K})^{(-)}
\ltimes_\phi T(r_\Phi)$, где $\phi$ "--- присоединенное действие. Остается
задать на идеале $T(r_\Phi)$ ассоциативное умножение $\cdot_\Phi$ с условием
$\bigl(T(r_\Phi),\cdot_\Phi\bigr)^{(-)}=T(r_\Phi)$, согласованное с $\phi$. Для
типов $D_n$ и $C_n$ идеал $T(r_\Phi)$ абелев, и на нем задаем нулевое
умножение; для типа $B_n$ коммутант $T(\epsilon_1)$ лежит в его центре, и
полагаем $e_r\cdot_{B_n}e_s=\frac12[e_r,e_s]$ при $\Char\mathbb{K}\neq2$ и
$e_r\cdot_{B_n}e_s=0$ иначе. Оператор $\phi_x$ при этом является
дифференцированием алгебры $\bigl(T(r_\Phi),\cdot_\Phi\bigr)$ по тождеству
Якоби.

\dissth{th:BnCnDn-pre-lie}

Построенная алгебра $R\Phi'(\mathbb{K})$ нильпотентна и градуирована решеткой
корней согласованно с градуировкой Картана; она неассоциативна для типов $C_n$
$(n\geq3)$, $D_n$ $(n\geq4)$ и $B_n$ $(n\geq3,\ \Char\mathbb{K}\neq2)$, что
видно из \crefRod{th:semidirect-assoc-cond}. Заметим, что левосимметрическая
структура на $N\Phi(\mathbb{K})$ получается и из невырожденного
дифференцирования $D(e_r)=\height(r)e_r$ по формуле $x\cdot y=D^{-1}[x,Dy]$, но
лишь при $\Char\mathbb{K}=0$ либо $\Char\mathbb{K}>\height(\rho)$; построенная
выше конструкция от ограничений на характеристику свободна.

\medskip

\textbf{Четвертая глава} посвящена комбинаторному перечислению идеалов.

В \S~4.1 подпространство $S$ в $\mathbb{K}^m$ названо \emph{координатно
    полным}, если оно не лежит ни в одной координатной гиперплоскости; число таких
$t$-мерных подпространств в $\mathbb{F}_q^m$ обозначено $V_q(m,t)$. Формулой
включений-исключений по координатным гиперплоскостям получено
\[
    V_q(m,t) = \sum_{k=t}^{m} (-1)^{m-k} \binom{m}{k} {k \brack t}_{q},
\]
что дает алгебраически-комбинаторное доказательство формул, полученных ранее
Г.\:П.~Егорычевым методом коэффициентов.

Всякий стандартный идеал алгебры Ли $N\Phi(\mathbb{K})$ с множеством углов
$\mathcal{L}$ порядка $m$ имеет вид $H(\mathcal{L},S)=Q(\mathcal{L})+\{
a_1e_{r_1}+\dots+a_me_{r_m} \mid (a_1,\dots,a_m)\in S\}$ для однозначно
определенного координатно полного подпространства $S$ в $\mathbb{K}^m$. Поэтому
в \S~4.2 перечисление сводится к числам $B(\Phi,m)$ $m$-элементных множеств
углов в $\Phi^+$ и числам $V_q(m,t)$, и проблема~1 из \cite{sigsam2001}
решается.

\dissth{th:enum-stand-id-classical}

\dissth{th:enum-stand-id-exceptional}

В \S~4.3 для произвольной ненулевой подалгебры алгебры Ли $N\Phi(\mathbb{K})$
строится специальная каноническая база: базу пересечения с $Q(\mathcal{L})$
дополняют элементами $\alpha_i=\sum_j a_{ij}e_{r_j}+\alpha_i'$, где
$\alpha_i'\in Q(\mathcal{L})$, а матрица $\|a_{ij}\|$ приведена к ступенчатому
виду.

Для формулировки результата о нестандартных идеалах точной обертывающей алгебры
$RD_n(\mathbb{K})$ используем реализацию
$D_n^+=\{\epsilon_a\pm\epsilon_b\mid1\leq b<a\leq n\}$. Черта над корнем
обозначает замену $\epsilon_1$ на $-\epsilon_1$ при фиксированных остальных
$\epsilon_b$. Рассмотрим множество углов
\begin{equation*}
    \mathcal{L}=\{r_1=s,r_2=\bar{s},r_3,\dots,r_m\}\quad(2\leq m<n),
    \tag{\ref*{eq:R-stepsId}}
\end{equation*}
где $s=\epsilon_a-\epsilon_1$ $(2\leq a<n)$ и $\bar{r}_j=r_j$ при
$3\leq j\leq m$. Порядок остальных углов фиксируем. Положим
$p=\epsilon_2-\epsilon_1$, $p_0=\epsilon_3-\epsilon_2$ и
$s_j=\epsilon_{a+j}-\epsilon_1$ $(1\leq j\leq n-a)$. При $s=p$ полагаем
$s_0=s+p_0$, а при $s\neq p$ "--- $s_0=s$.

Для нестандартного идеала $H$ с множеством углов $\mathcal{L}$ выбираем
наименьший номер $k\geq1$ с условием $Q(s_k)\subset H$ и полагаем
$t=\dim_{\mathbb{K}}\bigl(H/(H\cap Q(\mathcal{L}))\bigr)$; при этом $1\leq
k\leq n-a$ и $1\leq t<m$. Матрица $\|a_{ij}\|$ над $\mathbb{K}$ имеет размер
$t\times m$, ранг $t$, приведенный ступенчатый вид и не имеет нулевых столбцов;
$a_{12}=c\neq0$, $c,d_1,\dots,d_t\in\mathbb{K}$. Для записи идеала используем
обозначение
\begin{equation*}
    \begin{gathered}
        Id\ \{\mathcal{L},k,t,\|a_{ij}\|,c,d_1,\dots,d_t\}\coloneq
        \sum_{i=1}^t\mathbb{K}\left(d_ie_{s_k}+\sum_{j=1}^m a_{ij}e_{r_j}\right)+{}\\
        {}+T(s_0+\bar{p})+Q(s_k)+Q(\bar{s}_k)
        +\sum_{j=1}^k\mathbb{K}(e_{s_j}+ce_{\bar{s}_j})+\sum_{j=3}^m Q(r_j).
    \end{gathered}
    \tag{\ref*{eq:IdH0}}
\end{equation*}

\dissth{th:ExIdForm}

В \S~4.4 описание нестандартных идеалов наборами параметров используется для их
перечисления над конечным полем. Возникающие комбинаторные суммы вычисляются
методом коэффициентов, что приводит к следующей теореме.

\dissth{th:Enum-Dn}

Вместе с \crefTvor{th:enum-stand-id-classical} это завершает для каждого
классического типа перечисление всех идеалов одной из специальных точных
обертывающих алгебр.

\medskip

\FloatBarrier

%%%%%%%%%%%%%%%%%%%%%%%%%%%%%%%%%%%%%%%%%%%%%%%%%%%%%%%%%%%%%%%%%%%%%%%%%%%%%%
% Списки литературы
%%%%%%%%%%%%%%%%%%%%%%%%%%%%%%%%%%%%%%%%%%%%%%%%%%%%%%%%%%%%%%%%%%%%%%%%%%%%%%
\pdfbookmark{Литература}{bibliography}

% Работы автора (refsection=0). Порядок \nocite задает порядок в списке.
\nocite{vmj2015}%
\nocite{jsfu2018}%
\nocite{imm2020}%
\nocite{smj2023}%
\nocite{jsfu2026}%
\nocite{svobodniy2016}%
\nocite{levconf2016}%
\nocite{malcevconf2016}%
\nocite{modernWorldConf2017}%
\nocite{kuroshconf2018}%
\nocite{malcevconf2023}%

\ifdefmacro{\microtypesetup}{\microtypesetup{protrusion=false}}{}
\urlstyle{rm}
\insertbiblioexternal   % процитированные в автореферате чужие работы
\insertbiblioauthor     % публикации автора по теме диссертации
\ifdefmacro{\microtypesetup}{\microtypesetup{protrusion=true}}{}
\urlstyle{tt}
